\Askhsh{31149_SOLUTION}{1}
\begin{ylh}{1}
    \item [Ύλη από εικόνα]
\end{ylh}
\textbf{ΛΥΣΗ}

\begin{alist}
    \item Παρατηρούμε ότι για $x > 0$ είναι $1 + \frac{1}{x} > 1 \Leftrightarrow \ln \left(1 + \frac{1}{x}\right) > \ln 1 = 0$, άρα $f(x) = \frac{\ln\left(1+\frac{1}{x}\right)}{x^2} > 0$.
    Για κάθε $x > 0$ έχουμε
    \begin{align*}
        f'(x) &= \left[\frac{\ln\left(1+\frac{1}{x}\right)}{x^2}\right]' \\
        &= \frac{\left[\ln\left(1+\frac{1}{x}\right)\right]' \cdot x^2 - \ln\left(1+\frac{1}{x}\right)\cdot (x^2)'}{(x^2)^2} \\
        &= \frac{\left(\frac{1}{1+\frac{1}{x}} \cdot \left(-\frac{1}{x^2}\right)\right) \cdot x^2 - \ln\left(1+\frac{1}{x}\right)\cdot 2x}{x^4} \\
        &= \frac{\left(\frac{x}{x+1} \cdot \left(-\frac{1}{x^2}\right)\right) \cdot x^2 - \ln\left(1+\frac{1}{x}\right)\cdot 2x}{x^4} \\
        &= \frac{-\frac{x}{x+1} - 2x \ln\left(1+\frac{1}{x}\right)}{x^4} \\
        &= -\frac{\frac{1}{x+1} + 2 \ln\left(1+\frac{1}{x}\right)}{x^3}
    \end{align*}
    Επειδή $\ln\left(1+\frac{1}{x}\right) = x^2 f(x)$, η παράσταση γίνεται:
    \[
        f'(x) = -\frac{\frac{1}{x+1} + 2x^2 f(x)}{x^3}
    \]
    Παρατηρούμε ότι ο αριθμητής $\left(\frac{1}{x+1} + 2x^2 f(x)\right)$ και ο παρονομαστής $x^3$ του κλάσματος είναι θετικοί αριθμοί για $x > 0$.
    Συνεπώς, $f'(x) < 0$, για κάθε $x \in (0, +\infty)$.
    Ώστε η $f$ είναι γνησίως φθίνουσα στο $(0, +\infty)$. \monades{X}

    \item Η ανίσωση γράφεται ισοδύναμα
    \[
        \ln \left(\frac{1+f(x)}{f(x)}\right) > f^2(x) \cdot f(\ln2) \Leftrightarrow \frac{\ln\left(\frac{1}{f(x)}+1\right)}{f^2(x)} > f(\ln2) \Leftrightarrow f(f(x)) > f(\ln2)
    \]
    και καθώς η $f$ είναι γνησίως φθίνουσα παίρνουμε $f(x) < \ln2$.
    Αλλά $f(1) = \frac{\ln\left(1+\frac{1}{1}\right)}{1^2} = \ln2$.
    Έτσι, η τελευταία ανίσωση γίνεται $f(x) < f(1) \Leftrightarrow x > 1$. \monades{X}

    \item Καθώς $f(x) > 0$ για κάθε $x > 0$, το ζητούμενο εμβαδόν θα είναι $E = \int_{1/2}^{1} f(x)dx$.
    Θέτουμε $u = 1 + \frac{1}{x}$. Τότε $du = -\frac{1}{x^2}dx$.
    Για $x = \frac{1}{2}$, παίρνουμε $u = 1 + \frac{1}{1/2} = 1+2 = 3$.
    Για $x = 1$, παίρνουμε $u = 1 + \frac{1}{1} = 2$.
    Έτσι, έχουμε:
    \begin{align*}
        E &= \int_{3}^{2} \ln u \cdot (-du) \\
        &= \int_{2}^{3} \ln u \cdot du
    \end{align*}
    Εφαρμόζουμε ολοκλήρωση κατά μέρη:
    \begin{align*}
        E &= \int_{2}^{3} (u)' \cdot \ln u \cdot du \\
        &= [u \cdot \ln u]_{2}^{3} - \int_{2}^{3} u \cdot (\ln u)' du \\
        &= (3\ln3 - 2\ln2) - \int_{2}^{3} u \cdot \frac{1}{u} du \\
        &= 3\ln3 - 2\ln2 - \int_{2}^{3} 1 \cdot du \\
        &= \ln(3^3) - \ln(2^2) - [u]_{2}^{3} \\
        &= \ln27 - \ln4 - (3-2) \\
        &= \ln \left(\frac{27}{4}\right) - 1 \\
        &= \ln \left(\frac{27}{4}\right) - \ln e \\
        &= \ln \left(\frac{27}{4e}\right)
    \end{align*} \monades{X}
\end{alist}